% Generato da scripts/build_thesis_tables.py (T65) dai CSV versionati: non modificare a mano.
% Richiede \usepackage{booktabs}.
\begin{table}[htbp]
\centering
\small
\caption{Modulo 2: decomposizione di $\delta$ attorno alla sospensione del 2019-20}
\label{tab:modulo2_covid}
\begin{tabular}{lrc}
\toprule
Quantità & Stima & IC 95\% \\
\midrule
totale & $-$0,0050 & [$-$0,0395; 0,0350] \\
compositiva & 0,0116 & [$-$0,0169; 0,0383] \\
strutturale & $-$0,0004 & [$-$0,0390; 0,0395] \\
interazione & $-$0,0163 & [$-$0,0502; 0,0188] \\
quota compositiva & 0,4115 & [0,0339; 0,5580] \\
quota strutturale $\pi$ & 0,0136 & [0,0228; 0,7961] \\
\bottomrule
\end{tabular}

\smallskip
\begin{minipage}{\linewidth}
\footnotesize
Scenario 2 di T48: A = 288 partite del 2019-20 fino al 2020-03-09; B = 472 partite dal 2020-06-17 (resto del 2019-20 e 2020-21).\par
Oaxaca-Blinder three-fold di $\delta = \ln \hat p - \ln q$ (q Shin): totale = compositiva + strutturale + interazione; $\pi = |S| / (|C| + |S| + |I|)$, quota compositiva $|C| / (|C| + |S| + |I|)$.\par
IC al 95\%: stationary bootstrap dentro ciascun gruppo, 999 ricampionamenti, percentili.\par
Fonte: \texttt{results/us\_c7\_3b\_known\_scenarios.csv} (SHA256 \texttt{4365678b0eb3}).
\end{minipage}
\end{table}
